\clearpage
\VerdierPhysicalPageAnchor{0155}%
\begingroup\small
\setlength{\parskip}{0pt}\setlength{\emergencystretch}{1.5em}

\VerdierSemanticAnchor{ch3-3-2-2-admissible}
Un objet de $\mathcal D^\wedge$ est dit \emph{admissible} s'il est image d'un
morphisme entre foncteurs représentables. Soit
$\operatorname{Im}\colon\mathcal{F}\ell_1(\mathcal D)\to\mathcal D^\wedge$
le foncteur
\VerdierSemanticAnchor{ch3-3-2-2-1}
\[
\operatorname{Im}\bigl(X\xrightarrow fY\bigr)
=\text{image du morphisme }h(f).
\tag{3.2.2.1}
\]
Le foncteur $\operatorname{Im}$ est additif et transforme les morphismes
négligeables en morphismes nuls. Il se factorise donc d'une manière unique par
un foncteur additif noté abusivement
\VerdierSemanticAnchor{ch3-3-2-2-2}
\[
\operatorname{Im}\colon\mathrm A(\mathcal D)\longrightarrow\mathcal D^\wedge.
\tag{3.2.2.2}
\]
Le foncteur composé
$\operatorname{Im}\circ\mathcal H\colon
\mathcal D\longrightarrow\mathrm A(\mathcal D)
\longrightarrow\mathcal D^\wedge$
n'est autre que le foncteur $h$. Un objet de
$\mathcal D^\wedge$ est isomorphe à l'image par $\operatorname{Im}$ d'un objet de
$\mathrm A(\mathcal D)$ si et seulement s'il est admissible.

\VerdierSemanticAnchor{ch3-3-2-3}
\medskip\noindent\textbf{Proposition 3.2.3.}\quad
\textit{Le foncteur
$\operatorname{Im}\colon\mathrm A(\mathcal D)\to\mathcal D^\wedge$ est
pleinement fidèle.}

\VerdierSemanticAnchor{ch3-3-2-3-morphism-square}
Soit
\[
\begin{tikzcd}[column sep=4em,row sep=2.4em]
X\arrow[r,"f"]\arrow[d,"g"']&Y\arrow[d,"g'"]\\
X'\arrow[r,"f'"']&Y'
\end{tikzcd}
\qquad=(g;g')
\]
un morphisme de $\mathcal{F}\ell_1(\mathcal D)$. Nous noterons
$\operatorname{Im}((g;g'))$ son image. Le diagramme

\VerdierSemanticAnchor{ch3-3-2-3-canonical-diagram}
\[
\begin{tikzcd}[column sep=4em,row sep=2.2em]
X\arrow[r,"\operatorname{id}_X"]\arrow[d,"\operatorname{id}_X"']&X\arrow[d,"f"]\\
X\arrow[r,"f"]\arrow[d,"f"']&Y\arrow[d,"\operatorname{id}_Y"]\\
Y\arrow[r,"\operatorname{id}_Y"']&Y
\end{tikzcd}
\]
fournit, après application de $\operatorname{Im}$,
\[
h(X)\xrightarrow{\operatorname{Im}((\operatorname{id}_X;f))}
\operatorname{Im}(f)\xrightarrow{\operatorname{Im}((f;\operatorname{id}_Y))}h(Y).
\]

\VerdierSemanticAnchor{ch3-3-2-3-1}
On a l'égalité
\[
\operatorname{Im}((f;\operatorname{id}_Y))\circ
\operatorname{Im}((\operatorname{id}_X;f))=h(f).
\tag{3.2.3.1}
\]
\endgroup
